\subsection{Notations et hypothèses}
\label{sec:notations}

Soit :
\begin{itemize}
    \item \(C_0\): capital initial investi;
    \item \(n\) : durée de détention (en années) ;
    \item \(r\): rendement brut annuel;
    \item \(f_g\) : frais de gestion annuels (assurance-vie) ;
    \item \(t_{ps}\) : taux de prélèvements sociaux ;
    \item \(A_s\): abattement successoral (CTO);
    \item \(A_{av}\): abattement assurance-vie (152 500 € par bénéficiaire).
\end{itemize}
Les valeurs finales des capitaux sont alors :
\begin{align}
    C_{CTO} &= C_0(1 + r)^n, \\
    C_{AV} &= C_0(1 + r - f_g)^n.
\end{align}
La plus-value en assurance-vie est donnée par \(\Delta = C_{AV} – C_0\), et les prélèvements sociaux par \(PS = t_{ps} \cdot \Delta^+\). Le capital net après prélèvements sociaux est \(X = C_{AV} - PS\).
